% Generated draft: requires mathematical and source review.
\documentclass[11pt]{article}
\usepackage[margin=25mm]{geometry}
\usepackage{fontspec}
\setmainfont{Latin Modern Roman}
\usepackage{amsmath,amssymb,mathtools,unicode-math}
\setmathfont{Latin Modern Math}
\usepackage{enumitem,needspace,xurl,hyperref}
\setlength{\parindent}{0pt}
\setlength{\parskip}{5pt}
\setlength{\emergencystretch}{3em}
\newcommand{\N}{\mathbb{N}}
\newcommand{\Z}{\mathbb{Z}}
\newcommand{\Q}{\mathbb{Q}}
\newcommand{\R}{\mathbb{R}}
\newcommand{\C}{\mathbb{C}}
\newcommand{\F}{\mathbb{F}}
\newcommand{\A}{\mathbb{A}}
\newcommand{\PP}{\mathbb{P}}
\newcommand{\E}{\mathbb{E}}
\newcommand{\eps}{\varepsilon}
\newcommand{\dd}{\,\mathrm{d}}
\begin{document}
\section*{UnsolvedMath ID: 30002482}
\textbf{Source problem code:} catalogue entry
\subsection{Definitions and mathematical statement}
For a representation \(\theta\) of a group \(G\) on a vector space \(V_\theta\) over a coefficient field \(k\), and a subgroup \(H\subseteq G\), an \(H\)-invariant linear functional is a map \(\ell:V_\theta\to k\) satisfying
\[
\ell(\theta(h)v)=\ell(v)\qquad(h\in H,\ v\in V_\theta).
\]
A period means a nonzero such functional, subject to any continuity requirements of the intended representation category. In the general linear group setting, the usual symplectic subgroup is
\[
\operatorname{Sp}_{2n}(F)=\{g\in\operatorname{GL}_{2n}(F):g^{\mathsf T}J_ng=J_n\},\qquad
J_n=\begin{pmatrix}0&I_n\\-I_n&0\end{pmatrix}.
\]
Thus, if this is the intended ambient setting, having a symplectic period means
\[
\operatorname{Hom}_{\operatorname{Sp}_{2n}(F)}(V_\theta,k)\ne\{0\},
\]
where \(k\) carries the trivial action and the Hom space is interpreted in the intended category. The supplied record does not specify \(F\), \(k\), the ambient group, or that category; this description is conditional and does not supply those missing hypotheses.

Retain the source notation \(Q(\Delta_1,\ldots,\Delta_{2r})\). For a nonempty list of length \(2r\), the adjacent-pair condition is
\[
Q(\Delta_{2i+1},\Delta_{2i+2})\text{ is essentially Speh}
\qquad\text{for every }i\in\{0,\ldots,r-1\}.
\]
Let \(\mathcal G'\) denote the source's class of representations \(Q(\Delta_1,\ldots,\Delta_{2r})\) built from ladder representations and satisfying this condition. The question is whether
\[
\forall\theta\in\mathcal G',\qquad
\operatorname{Hom}_{H_{\mathrm{sp}}}(V_\theta,k)\ne\{0\},
\]
where \(H_{\mathrm{sp}}\) denotes the intended symplectic subgroup of the ambient group of \(\theta\), and \(k\) is the intended coefficient field with trivial action.

The supplied record does not define the objects \(\Delta_j\), the construction \(Q\), ladder representations, or the qualification ``essentially Speh.'' Consequently, it does not determine membership in \(\mathcal G'\) precisely. It also does not specify whether \(r=0\) is allowed. These qualifications remain unresolved; the displayed assertion records the source's intended question rather than a complete standalone conjecture with fixed ambient data. The dated literature assessment additionally describes the representations as ``generic,'' but the supplied mathematical statement does not explicitly impose genericity.
\subsection{Short English statement}
The question asks whether every representation in the specified ladder-based class, whose adjacent pairs are essentially Speh, admits a nonzero invariant functional for the relevant symplectic subgroup. The ambient representation setting and several defining terms are missing, so this remains a qualified statement-only draft requiring mathematical review.
\subsection{Sources}
\begin{itemize}
\item Supplied local registry record: permanent UnsolvedMath ID \(30002482\), source problem code \texttt{catalogue entry}, title ``Symplectic Periods of Essentially Speh Representations.'' The record gives proposed year 2014 and names no proposer.
\item Supplied report citation: \emph{Representation Theory and Analysis of Reductive Groups: Spherical Spaces and Hecke Algebras} (2014), Oberwolfach Reports, pp.\ 145--212, DOI \texttt{10.4171/owr/2014/03}. \url{https://doi.org/10.4171/owr/2014/03}. The background describes Oberwolfach Reports as documenting mathematical workshops and their open-problem discussions.
\item Supplied literature source, ``Oberwolfach representation theory report'': \url{https://ems.press/content/serial-article-files/46494?nt=1}.
\item Supplied literature source, ``Mitra, symplectic periods for GL representations'': \url{https://arxiv.org/abs/1601.03611}.
\item Literature assessment dated exactly \texttt{2026-08-21}: the supplied curation labels the problem open and says that the symplectic-period assertion for the stated generic representations is posed as a conjecture beyond classified low-rank cases, with no general proof found. This is supplied dated curation derived from the OWR material, not an independently verified literature assessment. The background expressly requests independent verification of time-sensitive status and cited claims.
\item The registry labels the statement \texttt{corrected\_verified} and reports manual semantic, typographic, and LaTeX review of a meaning-preserving cleanup. These are supplied editorial qualifications, not independent verification of mathematical completeness or correctness.
\end{itemize}
\end{document}
